% SPDX-License-Identifier: CC-BY-4.0
% =====================================================================
% Glossary for D. Joyner and Claude, "A hitch-hiker's guide to the
% metaplectic trace formula for SL(2)" (November 2026).
% Added in response to the second referee report (ChatGPT). Needs longtable.
% Background pass of November 18, 2026: the list of terms was rewritten and
% enlarged. Each entry is a short summary that can be read on its own, with
% no pointers to numbered statements. The table of notation keeps its
% "Defined in" column.
% =====================================================================

\section*{Glossary}\label{sec:glossary}
\addcontentsline{toc}{section}{Glossary}

\subsection*{Notation}

{\small
\begin{longtable}{@{}p{0.21\textwidth}p{0.55\textwidth}p{0.18\textwidth}@{}}
Symbol & Meaning & Defined in\\\hline
\endhead
\multicolumn{3}{@{}l}{\emph{Groups}}\\
$k$, $d$, $\A_k$, $\A$ & number field, its degree $[k:\Q]$, its adeles,
  the adeles of $\Q$ & \S\ref{sec:background}\\
$G$ & $\Res_{k/\Q}SL(2)$; $G(\Q)=SL(2,k)$ & \S\ref{sec:background}\\
$B$, $N$, $A$ & upper triangular Borel subgroup, its unipotent radical,
  the diagonal torus & \S\ref{sec:local}\\
$a(t)$, $n(u)$, $w$ & $\diag(t,t^{-1})$, the upper unipotent matrix with
  entry $u$, the Weyl element $\bigl(\begin{smallmatrix}0&-1\\1&0
  \end{smallmatrix}\bigr)$ & \S\ref{sec:local}, \S\ref{sec:geometric}\\
$T_\theta$, $T_\mu$ & anisotropic torus attached to $k(\sqrt\theta)$;
  centralizer of an elliptic $\mu$ & \S\ref{sec:local}, Lemma~\ref{lem:classes}\\
$K_v$, $K_m$, $K_0$, $K_f$ & maximal compact subgroup at $v$; principal
  congruence subgroup of level $\varpi^m$; $\prod_vK_v$; an open compact
  subgroup of $G(\A_f)$ & \S\ref{sec:local}, \S\ref{sec:global}\\
$A(\A)^1$, $A_B(\R)^0$ & kernel of $H$ on $A(\A)$; the split component
  $\{a(u)\mid u>0\}$ & \S\ref{sec:torus-cover}, \S\ref{sec:coarse}\\[4pt]
\multicolumn{3}{@{}l}{\emph{The cover}}\\
$\mu_2$ & $\{\pm1\}$, the kernel of the covering & \S\ref{sec:background}\\
$(a,b)_v$ & quadratic Hilbert symbol of $k_v$ & \S\ref{sec:background}\\
$\beta_v$ & Kubota's cocycle on $SL(2,k_v)$; Borel, not continuous &
  \eqref{eq:beta}, Lemma~\ref{lem:kubota}\\
$\kappa_v$, $\kappa_\A$ & Kubota's correction on $K_v$ ($v$ finite,
  non-dyadic); its product over $v$, defined on $G(\Q)$ &
  \eqref{eq:kappa}, \eqref{eq:kappa-A}\\
$\sigma_v$, $\sigma_\A$ & $\beta_v$ times the coboundary of $\kappa_v$;
  the adelic cocycle $\prod_v\sigma_v$ & \eqref{eq:sigma-v},
  \eqref{eq:sigma-A}\\
$\overline{SL(2,k_v)}$, $\overline{G(\A)}$ & the local and the adelic
  double cover; $\overline H=p^{-1}(H)$ for a subgroup $H$ &
  Lemma~\ref{lem:kubota}, \S\ref{sec:global}\\
$(g,\zeta)$ & Borel coordinates on a cover; $g\mapsto(g,1)$ is not
  continuous & \eqref{eq:law}, Example~\ref{ex:discontinuous}\\
$K_v^*$, $K_f^*$ & the image of $K_v$ under $k\mapsto(k,\kappa_v(k))$, a
  compact open subgroup of the cover; the image of $K_f$ &
  Lemma~\ref{lem:local}(b), Lemma~\ref{lem:global}(c)\\
$\overline{G(\Q_S)}$, $\overline{A(\Q_S)}$ & the covers of
  $\prod_{v\in S_k}SL(2,k_v)$ and of its diagonal torus &
  \S\ref{sec:local-dist}, \S\ref{sec:invariant}\\[4pt]
\multicolumn{3}{@{}l}{\emph{Heights, measures, functions}}\\
$H=H_B$, $H_v$, $H_S$ & the height $\log|t|$ of the $A$-component in
  $x=nak$ & \eqref{eq:H}, Lemma~\ref{lem:NAK-global}\\
$\delta_B$ & modulus character, $e^{2H}$ & \S\ref{sec:measures-local}\\
$\mathfrak a$, $\alpha$, $\alpha^\vee$ & $\Hom(X(A)_\Q,\R)\cong\R$; the
  simple root $2/d$; the coroot $d$ & \S\ref{sec:chambers},
  Remark~\ref{rem:alpha}\\
$\tau_B=\hat\tau_B$ & characteristic function of $(0,\infty)$ &
  \eqref{eq:tau}\\
$C_c^\infty(\cdot)_-$ & smooth, compactly supported, anti-genuine
  functions on a cover & \S\ref{sec:functions}\\
$e_v$, $e^S$ & unit of the spherical Hecke algebra at $v$;
  $\bigotimes_{v\notin S}e_v$ & \eqref{eq:unit}, \S\ref{sec:local-dist}\\
$\mathcal H(\overline{G(\A)})$, $\mathcal H(\overline{G(\Q_S)})$ & Hecke
  algebras of test functions; $f=f_S\otimes e^S$ & \S\ref{sec:global},
  \S\ref{sec:local-dist}\\
$f^y$ & $x\mapsto f(yxy^{-1})$ & \S\ref{sec:invariant}\\
$c_k$ & $\vol(A(\Q)\backslash A(\A)^1)$ & \S\ref{sec:coarse}\\[4pt]
\multicolumn{3}{@{}l}{\emph{Representations and classical forms}}\\
$\Theta_\pi$, $\omega_\pi$ & the character $f\mapsto\operatorname{tr}\pi(f)$
  of a representation; its central character & \S\ref{sec:rep-local}\\
$\pi_\lambda$, $I(\pi_\lambda)$, $\pi^w$ & $\pi\,e^{\lambda H}$ for a genuine
  character $\pi$ of the cover of $A$; the principal series; the Weyl
  conjugate of $\pi$ & \S\ref{sec:rep-local}\\
$M_v(\pi,\lambda)$, $\rho(\pi,\lambda)$ & local intertwining operator
  $I(\pi_\lambda)\to I(\pi^w_{-\lambda})$; the scalar
  $M_v(\pi^w,-\lambda)M_v(\pi,\lambda)$ & \S\ref{sec:rep-local}\\
$D_m^\pm$ & representation of the real cover of lowest weight $m$
  (highest weight $-m$), $m\in\frac12+\Z$, $m>0$ & \S\ref{sec:rep-local}\\
$\tilde k_\theta$, $\tilde g_z$ & the one-parameter subgroup above the
  rotations; the lift $(n(u)a(y^{1/2}),1)$ of a point $z=u+iy$ &
  \S\ref{sec:rep-local}, \S\ref{sec:rep-global}\\
$\theta(z)$, $j(\gamma,z)$ & $\sum_ne^{2\pi in^2z}$; Shimura's factor
  $\theta(\gamma z)/\theta(z)$ on $\Gamma_0(4)$ & \eqref{eq:shimura-j}\\
$S_{\kappa/2}(\Gamma)$, $T(p^2)$ & cusp forms of weight $\kappa/2$ in the
  sense of Shimura; Shimura's Hecke operator & \S\ref{sec:rep-global}\\
$A_t(n)$, $F_t$ & the coefficients of the Shimura lift of a form with
  coefficients $a(n)$; the lift & \eqref{eq:shimura-lift}\\[4pt]
\multicolumn{3}{@{}l}{\emph{Rank $0$}}\\
$\psi$, $\psi_b$ & a nontrivial character of $\A_k/k$; $x\mapsto\psi(bx)$
  & \S\ref{sec:torus-cover}\\
$\gamma_v(a)$ & Weil index of $x\mapsto\psi_v(ax^2)$, an eighth root of
  unity & \eqref{eq:weil-index}\\
$\chi_\psi$, $\chi_{\psi_b}$ & Weil's genuine character of
  $\overline{A(\A)}$, $(a(t),\zeta)\mapsto\zeta\prod_v\gamma_v(t_v)/
  \gamma_v(1)$ & Lemma~\ref{lem:weil}\\
$X_T$, $\Pi_{\mathrm{aut}}(X_T)$ & $T(\Q)\backslash\overline{T(\A)^1}$;
  its genuine characters & \S\ref{sec:torus-cover}\\
$I^A_{\mathcal O}$, $I^A_{\bar\chi}$ & the two kinds of terms in the
  trace formula for the cover of $A$ & \eqref{eq:IA}\\[4pt]
\multicolumn{3}{@{}l}{\emph{Geometric side}}\\
$\mathcal C$, $\mathcal O$, $\mathcal O_z$ & the set of classes in
  $G(\Q)$; a class; the central class of $z=\pm1$ &
  Lemma~\ref{lem:classes}\\
$K_P$, $K_{P,\mathcal O}$, $J_{P,\mathcal O}$ & kernels attached to
  $P\in\{B,G\}$ & \eqref{eq:KP}, \eqref{eq:JP}\\
$k^T$, $k^T_{\mathcal O}$, $j^T_{\mathcal O}$ & truncated kernels &
  \eqref{eq:kT}\\
$J^T_{\mathcal O}(f)$, $J_{\mathcal O}(f)$ & integral of
  $k^T_{\mathcal O}$, a polynomial in $T$; its value at $T=0$ &
  Theorem~\ref{thm:coarse}\\
$v_A$, $v_{A,v}$ & the weight $-H(x)-H(wx)$ & \eqref{eq:vA},
  Lemma~\ref{lem:weight-local}\\
$\Phi_f$, $Z(f,s)$ & the function $y\mapsto\int_Kf(k^{-1}n(y)k)\,dk$ and
  its zeta integral (not to be confused with $\Phi_A$) & \eqref{eq:Phi},
  \eqref{eq:Zf}\\
$J_G(\gamma,f_S)$ & orbital integral on $\overline{G(\Q_S)}$ &
  \eqref{eq:JG}\\
$J_A(\gamma,f_S)$ & weighted orbital integral, weight $v_A$; not
  invariant & \eqref{eq:JA}\\
$J_A(z,f_S)$ & its regularized value at $z=\pm1$ &
  Lemma~\ref{lem:JA-central}\\
$u_\xi$, $a^M(S,\gamma)$ & the unipotent element $n(\xi)$; the global
  coefficients of the fine expansion & Theorem~\ref{thm:fine}\\[4pt]
\multicolumn{3}{@{}l}{\emph{Spectral side}}\\
$\mathcal X$, $\mathcal X_G$, $\mathcal X_A$ & classes of cuspidal data;
  those of $(G,\pi)$; those of $(A,\mu)$ & \S\ref{sec:decomp}\\
$\mu$, $\mu^w$, $\mu_\lambda$ & a genuine automorphic character of the
  cover of $A$; its Weyl conjugate; $\mu\,e^{\lambda H}$ &
  \S\ref{sec:decomp}\\
$\chi_b$ & the ramified class of $(A,\chi_{\psi_b})$ &
  Lemma~\ref{lem:w-action}\\
$\mathcal H_B(\mu)$, $I(\mu_\lambda)$ & the space of the principal series
  and the induced representation & \eqref{eq:HBmu}\\
$E(x,\varphi,\lambda)$ & Eisenstein series & \eqref{eq:Eis}\\
$M(w,\lambda)$, $M_v$ & global and local intertwining operators &
  \eqref{eq:Mw}, Lemma~\ref{lem:intertwining}\\
$m^S(\mu,\lambda)$ & $L^S(2\lambda,\chi^2)/L^S(2\lambda+1,\chi^2)$ &
  \eqref{eq:mS}\\
$\omega_{\psi'}$, $r^\pm_{\psi'}$ & Weil representation and its even and
  odd parts & \eqref{eq:weil-rep}\\
$V_b$ & residual representation of the class $\chi_b$, isomorphic to
  $r^+_{\psi_{-b}}$ & Theorem~\ref{thm:residual}\\
$\Lambda^T$ & truncation operator & \eqref{eq:trunc}\\
$K_{P,\chi}$, $k^T_\chi$, $J^T_\chi$, $J_\chi$ & spectral kernels, their
  truncation and its integral & \S\ref{sec:Jchi}\\
$M_B^T(\mu,\lambda)$ & $2T-M(w,\lambda)^{-1}M'(w,\lambda)$ &
  \eqref{eq:MT}\\[4pt]
\multicolumn{3}{@{}l}{\emph{Invariant form}}\\
$u'_B(k,y)$, $f_{B,y}$ & the weight $-H(ky)$; the function on the cover
  of $A$ that measures $J(f^y)-J(f)$ & \eqref{eq:uprime}, \eqref{eq:fBy}\\
$V(\pi)$, $I(\pi_\lambda)$ & local principal series of
  $\overline{G(\Q_S)}$ & \S\ref{sec:invariant}\\
$r_v(\pi_v,\lambda)$, $R_v$, $R_S$ & local normalizing factor; normalized
  operator $r_v^{-1}M_v$; $\bigotimes_{v\in S_k}R_v$ &
  Theorem~\ref{thm:normalization}\\
$S_{\mathrm{bad}}(\mu)$ & the places where $r_v$ is not explicit in
  Section~\ref{sec:invariant}; made explicit in Proposition~\normref{}(c) &
  \S\ref{sec:normalized}\\
$\rho(\pi,\lambda)$, $\omega(\pi)$, $h_\pi$ & the scalar
  $M_v(\pi^w,-\lambda)M_v(\pi,\lambda)$; the value $\pi((w,1)^2)\in\mu_4$;
  the factor of $|\rho|$ without zeros or poles in $|q^{-\lambda}|<1$ &
  Proposition~\normref{}\\
$m(\mu,\lambda)$ & global normalizing factor, $M(w,\lambda)=m\cdot R_S$ &
  \eqref{eq:global-m}\\
$\varphi_A(f,\pi)$ & $-\operatorname{tr}(R_S^{-1}R_S'\,I(\pi,f))$, a
  function of the \emph{character} $\pi$ of $\overline{A(\Q_S)}$ &
  \eqref{eq:phiA}\\
$\Phi_A(\tilde a,f)$ & the inverse Fourier transform of
  $\varphi_A(f,\cdot)$, a function of the \emph{element} $\tilde a$ of
  $\overline{A(\Q_S)}$ & \eqref{eq:PhiA}\\
$I_A(\gamma,f)$, $I_A(z,f)$ & $J_A-\Phi_A$, the invariant distributions &
  \eqref{eq:IAgamma}\\
$I_{\mathcal O}(f)$, $I_\chi(f)$ & the terms of the invariant trace
  formula & \eqref{eq:Iochi}\\[3pt]
\multicolumn{3}{@{}l}{\emph{Section~\ref{sec:stable}}}\\
$G_1$, $\delta_a$, $\Ad(h)$ & $GL(2)$; $\diag(a,1)$; the
  automorphism of the cover lifting conjugation by $h\in G_1$ &
  Lemma~\ref{lem:gl2}, Lemma~\ref{lem:Ad}\\
$f^h$, $f^{[h]}$, $f_z$, $f_-$ & $f\circ\Ad(h)$; $f(h\star\cdot)$;
  $f(z\,\cdot)$ for central $z$; $f_{-1}$ & \S\ref{sec:stable},
  \S\ref{sec:stable-conj}\\
$D^+(\gamma)$, $\eta_\gamma$ & $\det(1+\gamma)=2+\operatorname{tr}\gamma$;
  the character $a\mapsto(a,D^+(\gamma))_v$ & \eqref{eq:eta}\\
$h\star x$ & $\eta_x(\det h)\,\Ad(h)x$, the action whose orbits are
  the stable classes & \eqref{eq:star}\\
$\Delta_\psi$ & genuine function with values in $\mu_4$, constant on
  stable classes; $|D^+|^{1/2}$ times the character of $r^+_\psi-r^-_\psi$ &
  Lemma~\ref{lem:Delta}, Prop.~\ref{prop:weil-char}\\
$\kappa_E$, $N E^\times$ & the quadratic character of $E/F$ and its
  kernel, the norms & \S\ref{sec:stable}\\
$J^{\mathrm{st}}_G$, $J^\kappa_G$ & stable and $\kappa$-orbital
  integrals & \eqref{eq:Jst}, \eqref{eq:Jst-global}\\
$S_{\mathrm{ell}}(f)$ & the stable elliptic regular distribution &
  Theorem~\ref{thm:elliptic-stable}\\
$H_1$, $\kappa_a$ & the height on $GL(2)$; the character
  $a(t)\mapsto(a,t)$ of the cover of $A$ & \S\ref{sec:hyperbolic-gl2}\\
$\mathfrak X_S$, $L^S(s,\kappa)$, $\rho_S$, $c_S$ & quadratic Hecke
  characters unramified outside $S$; their partial $L$-functions; the
  residue and constant term of $\zeta_k^S$ at $1$ & \S\ref{sec:unipotent}\\
$J_G(zu_\xi,f_S)$, $J^\kappa_{\mathrm{unip}}(z,f_S)$ & normalized unipotent
  orbital integral; its $\kappa$-sum over $\xi$ & \eqref{eq:JG-unip},
  Prop.~\ref{prop:unipotent}\\
$U_h$, $\Theta_b$, $\Theta^{\mathrm{st}}_b$ & the action of $G_1(k)$ on
  automorphic forms; the span of the theta functions of $\psi_{-b}$; its
  stable character & \S\ref{sec:spectral-gl2}\\
\end{longtable}
}

\subsection*{Terms}

Each entry is a short reminder of what the term means in this paper. The
entries can be read independently of the text and of each other.

\begin{description}
\item[Admissible representation.] A representation of a $p$-adic group or
  of its cover in which every vector is fixed by some open subgroup, and
  the vectors fixed by any given compact open subgroup form a space of
  finite dimension. For such a representation $\pi$ and a test function
  $f$, the operator $\pi(f)$ has finite rank, so its trace is defined.
\item[Anti-genuine function.] A function $f$ on a cover with
  $f(x\zeta)=\zeta^{-1}f(x)$ for $\zeta\in\mu_2$. All test functions are
  anti-genuine. The inverse is there so that $\pi(f)=\int f(x)\pi(x)\,dx$
  does not vanish on genuine representations. For the double cover
  $\zeta^{-1}=\zeta$, so the condition is the same as \emph{genuine},
  $f(x\zeta)=\zeta f(x)$.
\item[Automorphic form (genuine).] A smooth function $\varphi$ on the
  adelic cover that is left invariant under the rational points $G(\Q)$,
  satisfies $\varphi(x\zeta)=\zeta\varphi(x)$ for $\zeta\in\mu_2$, is
  finite under the maximal compact subgroup and under the center of the
  enveloping algebra, and has moderate growth. For $k=\Q$ these functions
  correspond to classical forms of half-integral weight on the upper half
  plane.
\item[Automorphic representation (genuine).] An irreducible representation
  of the adelic cover that occurs in the space of genuine automorphic
  forms. It is a restricted tensor product of irreducible genuine
  representations of the local covers, almost all of them unramified.
\item[Borel coordinates.] The coordinates $(g,\zeta)$, with $g$ in the
  linear group and $\zeta\in\mu_2$, in which a cover is described by a
  cocycle. They give the group law and the measurable structure. They do
  not give the topology, because the cocycle is not continuous: a
  sequence $g_j\to g$ need not satisfy $(g_j,1)\to(g,1)$.
\item[Central character.] The character by which the center of a group
  acts in an irreducible representation. The center of the local cover
  of $SL(2)$ is the preimage of $\{\pm1\}$, a group of order $4$. It is
  cyclic at the real places and at some dyadic places, and there the
  central character of a genuine representation takes the values $\pm i$
  on the lifts of $-1$.
\item[Character of a representation.] The linear form $f\mapsto
  \operatorname{tr}\pi(f)$ on test functions. It is invariant under
  conjugation and determines an irreducible representation. The spectral
  side of the trace formula is a sum and integral of characters.
\item[Class (geometric).] A set of elements of $G(\Q)$ whose semisimple
  parts are conjugate. For $SL(2)$ there are three kinds. A central class
  consists of $z=\pm1$ times unipotent elements. A hyperbolic class is
  the conjugacy class of a diagonal element $a(t)$ with $t\ne\pm1$. An
  elliptic class is the conjugacy class of an element whose eigenvalues
  are not in $k$. A class is \emph{unramified} if it consists of
  semisimple elements.
\item[Coarse expansion.] The first form of each side of the trace
  formula: a sum over geometric classes $\mathcal O$ of distributions
  $J_{\mathcal O}(f)$, or over classes $\chi$ of cuspidal data of
  distributions $J_\chi(f)$. The terms are defined by truncation and are
  not yet written as orbital integrals or characters.
\item[Cocycle.] A function $\beta$ of two group elements with values in
  $\mu_2$ such that $(g_1,\zeta_1)(g_2,\zeta_2)=(g_1g_2,\beta(g_1,g_2)
  \zeta_1\zeta_2)$ is a group law. Kubota's cocycle on $SL(2,F)$ is a
  Hilbert symbol of the lower left entries of $g_1$, $g_2$ and $g_1g_2$,
  or of the lower right entries when the lower left ones vanish.
\item[Constant term.] The integral $\int_{N(\Q)\backslash N(\A)}
  \varphi(nx)\,dn$ of an automorphic form over the unipotent radical of
  the Borel subgroup. It describes $\varphi$ near the cusp.
\item[Cover, double cover.] A group $\overline G$ with a surjective
  homomorphism $p\colon\overline G\to G$ whose kernel is the central
  subgroup $\mu_2=\{\pm1\}$. The double cover of $SL(2,F)$ used here does
  not split: it is not $SL(2,F)\times\mu_2$, unless $F=\C$. It splits over
  unipotent subgroups, over small compact open subgroups, and, in the
  adelic case, over the rational points.
\item[Cusp form.] An automorphic form whose constant term vanishes. The
  cusp forms span the cuspidal part of the discrete spectrum.
\item[Cuspidal datum.] A pair $(G,\pi)$ with $\pi$ a cuspidal
  automorphic representation, or a pair $(A,\mu)$ with $\mu$ a genuine
  automorphic character of the cover of the diagonal torus $A$. The data
  $(A,\mu)$ and $(A,\mu^w)$, with $\mu^w$ the conjugate of $\mu$ by the
  Weyl element, are equivalent. The class of $(A,\mu)$ is \emph{ramified}
  if $\mu^w=\mu$. This happens exactly when $\mu$ is a Weil character
  $\chi_{\psi_b}$.
\item[Discrete series.] An irreducible unitary representation whose
  matrix coefficients are square integrable. For the real cover these are
  the representations of lowest weight $m$ or highest weight $-m$ with
  $m\in\{\frac32,\frac52,\dots\}$. They correspond to holomorphic and
  antiholomorphic forms of weight $m$.
\item[Dual group.] The complex group that parametrizes unramified
  representations. For the double cover of $SL(2)$ it is $SL(2,\C)$, the
  dual group of $PGL(2)=SO(3)$. It is not $PGL(2,\C)$, the dual group of
  $SL(2)$.
\item[Dyadic place.] A place of $k$ above $2$. The explicit splitting
  $\kappa_v$ of the cover over the maximal compact subgroup is not
  available there, and one uses a splitting over a smaller congruence
  subgroup. The unramified computation of the normalizing factors does
  not apply at a dyadic place either. There the explicit form of the
  factors is taken from the work of Szpruch on local coefficients.
\item[Eisenstein series.] The function $E(x,\varphi,\lambda)=\sum_\delta
  \varphi_\lambda(\delta x)$, summed over $B(\Q)\backslash G(\Q)$, built
  from a function $\varphi$ in a principal series. It converges for
  $\mathrm{Re}\,\lambda>1$ and continues meromorphically. The Eisenstein
  series on the line $\mathrm{Re}\,\lambda=0$ describe the continuous
  spectrum, and their residues at $\lambda=\frac12$ are theta functions.
\item[Elliptic element.] An element of $SL(2,F)$ whose eigenvalues do not
  lie in $F$. Its centralizer is a torus $T_\theta$ isomorphic to the
  group of elements of norm $1$ in a quadratic extension of $F$.
\item[Endoscopic group.] A smaller group whose stable orbital integrals
  account for the unstable part of the trace formula. For the double
  cover of $SL(2)$ there are two elliptic endoscopic data, attached to
  the central elements $\pm1$ of the dual group $SL(2,\C)$. Both have the
  group $SO(3)$ in its split form $PGL(2)$. They are discussed in this
  paper but not used.
\item[Fine expansion.] The second form of the geometric side, due to
  Arthur: each class contributes a finite sum of weighted orbital
  integrals on the group over $\Q_S$, multiplied by global coefficients
  $a^M(S,\gamma)$.
\item[$(G,S)$-equivalence.] Arthur's equivalence relation on $G(\Q)$, for
  a finite set $S$ of places. Two elements are equivalent if their
  semisimple parts are conjugate and their unipotent parts are conjugate
  in the centralizer over $\Q_S$. The equivalence classes index the terms
  of the fine expansion.
\item[Genuine character.] A character of the cover of an abelian group
  that is the identity on $\mu_2$. The genuine characters of the cover of
  the diagonal torus are the products of one of them, for example Weil's
  character, with the ordinary characters of the torus.
\item[Genuine representation.] A representation of a cover in which
  $\mu_2$ acts by $\zeta\mapsto\zeta$. A representation that is not
  genuine is trivial on $\mu_2$ and is a representation of the linear
  group.
\item[Half-integral weight (in the sense of Shimura).] A holomorphic
  function $F$ on the upper half plane has weight $\kappa/2$, with
  $\kappa$ odd, on a subgroup $\Gamma$ of $\Gamma_0(4)$ if $F(\gamma z)=
  j(\gamma,z)^\kappa F(z)$ for $\gamma\in\Gamma$. Here $j(\gamma,z)=
  \theta(\gamma z)/\theta(z)$ is the factor of the theta function
  $\theta(z)=\sum_ne^{2\pi in^2z}$. It is a square root of $cz+d$ times a
  Jacobi symbol and a fourth root of unity. Such forms are the genuine
  automorphic forms on the adelic double cover of $SL(2)$ over $\Q$ that
  have weight $\kappa/2$ at the real place.
\item[Hecke algebra.] An algebra of test functions under convolution.
  The \emph{spherical} Hecke algebra at an odd place consists of the
  anti-genuine functions that are invariant on both sides under the
  lifted maximal compact subgroup $K_v^*$. It is commutative, and its
  unit $e_v$ acts on a genuine representation as the projection onto the
  $K_v^*$-fixed vectors. In classical terms it is generated by the Hecke
  operator $T(p^2)$.
\item[Height.] The function $H(x)=\log|t|$, where $a(t)$ is the diagonal
  component of $x$ in the decomposition $x=nak$. It is invariant on the
  left under $N(\A)$ and under $B(\Q)$, and on the right under the
  maximal compact subgroup. On the upper half plane it is half the
  logarithm of the imaginary part.
\item[Hilbert symbol.] The symbol $(a,b)_v=\pm1$ on $k_v^\times\times
  k_v^\times$, equal to $1$ exactly when $ax^2+by^2=1$ has a solution in
  $k_v$. It is bimultiplicative and symmetric, it depends on $a$ and $b$
  only modulo squares, and the product of $(a,b)_v$ over all places is
  $1$ for $a,b\in k^\times$. The cover is built from it.
\item[Hyperbolic element.] An element of $SL(2,F)$ conjugate to a
  diagonal element $a(t)$ with $t\ne\pm1$. Its centralizer is the
  diagonal torus.
\item[Intertwining operator.] The operator $M(w,\lambda)$, or locally
  $M_v(\pi,\lambda)$, given by integration over the unipotent group
  after translation by the Weyl element. It maps the principal series of
  $\pi_\lambda$ to that of $\pi^w_{-\lambda}$ and commutes with the action
  of the group. The constant term of an Eisenstein series is
  $\varphi_\lambda+(M(w,\lambda)\varphi)_{-\lambda}$.
\item[Invariant distribution.] A linear form $I$ on test functions with
  $I(f^y)=I(f)$ for all $y$ in the group, where $f^y(x)=f(yxy^{-1})$.
  Orbital integrals and characters are invariant. Weighted orbital
  integrals are not.
\item[Invariant trace formula.] The form of the trace formula in which
  every term is an invariant distribution. It is obtained from the
  coarse and fine expansions by subtracting from each weighted orbital
  integral a correction built from the logarithmic derivative of the
  normalized intertwining operators.
\item[Iwasawa decomposition.] The decomposition $x=nak$ of an element of
  the group or of its cover, with $n$ upper triangular unipotent, $a$
  diagonal and $k$ in the maximal compact subgroup. The component $a$ is
  determined up to the compact part of the torus.
\item[Metaplectic group.] The nontrivial double cover of $SL(2,F)$, or
  of the symplectic group $Sp(2n,F)$, and the corresponding adelic cover.
  Over $\R$ it is the connected Lie group $Mp(2,\R)$. In the 1995 notes
  the term also covers the $n$-fold covers of $SL(2)$.
\item[Normalizing factor.] A scalar function $r_v(\pi_v,\lambda)$ such
  that the normalized operator $r_v^{-1}M_v$ is unitary and holomorphic
  on the imaginary axis, satisfies a functional equation, is rational on
  each type of the maximal compact subgroup, and fixes the unramified
  vector at almost all places. Three questions about it are separate:
  whether it exists, what its explicit formula is, and how the product of
  the local factors enters the global function $m(\mu,\lambda)$. For the
  double cover the factor is a constant times $L_v(2\lambda,\chi_v^2)/
  L_v(2\lambda+1,\chi_v^2)$.
\item[Orbital integral.] The integral $J_G(\gamma,f)$ of a test function
  over the conjugacy class of $\gamma$, normalized by a power of the
  discriminant. It is an invariant distribution. On the cover the
  orbital integral at a lift of $\gamma$ is the negative of the one at
  the other lift.
\item[Packet.] A finite set of irreducible representations that the
  stable trace formula does not separate. For the cover the packets are
  those of Waldspurger. The even and odd Weil representations
  $r^\pm_{\psi'}$ attached to one character form the simplest example.
\item[Principal series.] The representations $I(\pi_\lambda)$ induced
  from a genuine character $\pi$ of the cover of the diagonal torus,
  twisted by $e^{\lambda H}$. They are unitary for unitary $\pi$ and
  imaginary $\lambda$, and on the double cover they are then irreducible.
\item[Residual spectrum.] The part of the discrete spectrum that is not
  cuspidal. For this group it is spanned by the residues of Eisenstein
  series at $\lambda=\frac12$. It is the sum of the even Weil
  representations $r^+_{\psi_b}$, one for each class $b$ in
  $k^\times/k^{\times2}$, and its elements are theta functions.
\item[Satake parameter.] The conjugacy class in the dual group
  $SL(2,\C)$ that determines an unramified representation. For the
  principal series of $\pi=\chi_\psi\cdot(\chi\circ p)$ and $\lambda$ it
  is the class of $\diag(\chi(\varpi)q^{-\lambda},\chi(\varpi)^{-1}
  q^{\lambda})$.
\item[Shimura correspondence.] Shimura's map from Hecke eigenforms of
  half-integral weight $\kappa/2$ to eigenforms of even weight
  $\kappa-1$, under which the eigenvalue of $T(p^2)$ becomes the
  eigenvalue of $T(p)$. In terms of representations it is Waldspurger's
  correspondence, defined by theta series, between genuine automorphic
  representations of the double cover of $SL(2)$ and automorphic
  representations of $PGL(2)$ and of its inner forms. It depends on the
  choice of an additive character.
\item[Splitting.] A homomorphism $s$ from a subgroup $H$ of the linear
  group to the cover with $p\circ s=\mathrm{id}$. The cover splits over
  $N$ by $n\mapsto(n,1)$, over the maximal compact subgroup at odd places
  by $k\mapsto(k,\kappa_v(k))$, and over $G(\Q)$ by $\gamma\mapsto(\gamma,
  \kappa_\A(\gamma))$. It does not split over $SL(2,F)$ or over the
  diagonal torus.
\item[Stable conjugacy.] On $SL(2,F)$: conjugacy under $GL(2,F)$. On the
  cover: the equivalence relation whose classes are the orbits of the
  twisted action $h\star x=(\det h,\det(1+\gamma))_v\,\Ad(h)x$ of
  $GL(2,F)$, where $\gamma$ is the image of $x$. A stable class lies over
  a stable class of $SL(2,F)$ and contains exactly one of the two lifts
  of each of its elements. Conjugation by $GL(2,F)$ without the twist can
  exchange the two lifts, so it gives a coarser relation.
\item[Stable distribution.] An invariant distribution given by a class
  function that is constant on stable classes, or a stable orbital
  integral, which is the sum of the orbital integrals over the conjugacy
  classes in a stable class. Equivalently, an invariant distribution
  that is also invariant under the twisted action of $GL(2)$.
\item[Stable(r) form.] The trace formula with its terms organized by
  stable classes, before any transfer to $SO(3)$. In it the regular
  elliptic terms are $S(f)+S(f_-)$ with $S$ stable and $f_-$ the
  translate of $f$ by the central element $-1$.
\item[Strong approximation.] The statement that $G(\Q)G(\R)$ is dense in
  $G(\A)$ for $G=SL(2)$. It identifies the adelic quotient, divided by a
  compact open subgroup, with the quotient of the real group by a
  congruence subgroup. The same holds for the cover.
\item[Tamagawa measure.] The Haar measure on $G(\A)$ defined by an
  invariant differential form over $\Q$ and the self-dual measures of an
  additive character. For $SL(2)$ it gives $G(\Q)\backslash G(\A)$ the
  volume $1$.
\item[Test function.] A smooth, compactly supported, anti-genuine
  function on a local cover, or a finite sum of products of such
  functions over all places, almost all of them equal to the unit of the
  spherical Hecke algebra. Smoothness refers to the topology of the
  cover, not to the coordinates $(g,\zeta)$.
\item[Theta function.] The automorphic form $\theta_\phi(x)=\sum_{\xi\in
  k}(\omega_\psi(x)\phi)(\xi)$ attached to a Schwartz function $\phi$ on
  the adeles through the Weil representation. For $k=\Q$ and a Gaussian
  it is $y^{1/4}\sum_ne^{2\pi in^2z}$, of weight $\frac12$.
\item[Trace formula.] The identity between the geometric and the spectral
  expansion of the trace of a test function acting on the discrete
  spectrum, corrected by the contribution of the continuous spectrum.
\item[Truncation.] Subtracting from a kernel, or from an automorphic
  form, its constant term in the region $\{H(\delta x)>T\}$ near the
  cusp. The truncated kernel is integrable over the diagonal, and its
  integral is a polynomial in $T$ of degree at most one.
\item[Unipotent element.] An element $n(u)$, or a conjugate of one. The
  unipotent classes of $SL(2,F)$ other than the identity are indexed by
  $F^\times/F^{\times2}$. Their contribution to the trace formula is
  computed with a zeta integral and involves values at $1$ of quadratic
  Hecke $L$-functions.
\item[Unramified representation.] An irreducible genuine representation
  of the cover at an odd place that has a nonzero vector fixed by the
  lifted maximal compact subgroup $K_v^*$. The fixed vector is unique up
  to a scalar.
\item[Weight (of a vector).] A vector in a representation of the real
  cover has weight $m$ if the one-parameter group above the rotations
  acts on it by $e^{im\theta}$. In a genuine representation the weights
  lie in $\frac12+\Z$.
\item[Weight (of an orbital integral).] The function $v_A(x)=-H(x)-H(wx)$
  that is inserted into the orbital integral of a hyperbolic element by
  truncation.
\item[Weighted orbital integral.] An orbital integral with the weight
  $v_A$ inserted. It arises from truncation and is not invariant: under
  conjugation of the test function it changes by an ordinary orbital
  integral on the cover of the torus.
\item[Weil character.] The genuine character $\chi_\psi(a(t),\zeta)=
  \zeta\,\gamma_\psi(t)/\gamma_\psi(1)$ of the cover of the diagonal
  torus, built from Weil indices. It depends on the additive character
  $\psi$, it is fixed by the Weyl group, and globally it is trivial on
  the rational points of the torus.
\item[Weil index.] The eighth root of unity $\gamma_v(a)$ attached to the
  character of second degree $x\mapsto\psi_v(ax^2)$. It is defined by
  $\int\psi_v(ax^2)\,dx=\gamma_v(a)|2a|_v^{-1/2}$ and satisfies
  $\gamma_v(ab)\gamma_v(1)=\gamma_v(a)\gamma_v(b)(a,b)_v$.
\item[Weil representation.] The genuine representation $\omega_\psi$ of
  the cover on the Schwartz functions on $k_v$ in which $n(u)$ acts by
  multiplication by $\psi(ux^2)$ and the Weyl element by a Fourier
  transform. It is the sum of the \emph{even} part $r^+_\psi$ and the
  \emph{odd} part $r^-_\psi$, on the even and the odd functions, and both
  are irreducible.
\item[Weyl element.] The matrix $w=\bigl(\begin{smallmatrix}0&-1\\1&0
  \end{smallmatrix}\bigr)$. Conjugation by it sends $a(t)$ to
  $a(t^{-1})$. On the cover of the torus this conjugation is not
  inversion.
\item[Whittaker functional.] A linear form $\ell$ on a representation
  with $\ell(\pi(n(u))\phi)=\psi'(u)\ell(\phi)$ for a nontrivial
  character $\psi'$. On the double cover it is unique up to a scalar
  when it exists. The Weil representations $r^\pm_\psi$ have such functionals
  only for the characters $\psi_a$ with $a$ a square.
\end{description}
